% Generated by tex/build_swebok_structure.py from tex/chapters/ch01/03_ソフトウェア要求の基礎.tex
\subsubsection{サービス品質制約}

\term{サービス品質制約}とは、特定の技術名を指定せず、ソフトウェアが達成すべき品質水準を表す要求である。サービス品質という表現は、利用者や運用者が受け取るサービスの品質として考えるとわかりやすい。

\par\smallskip\noindent\refstepcounter{table}\label{tab:ch01-05}\textbf{表 \thetable: 第01章 ソフトウェア要求 表05}\par\smallskip
\begin{tabularx}{0.97\linewidth}{@{}p{.26\linewidth}X@{}}
\toprule
品質水準 & 要求文の例 \\
\midrule
応答時間 & 通常時、検索結果を2秒以内に返す。 \\
スループット & 1分あたり1万件のイベントを処理する。 \\
正確性 & 金額計算は定められた丸め規則に従う。 \\
信頼性 & 月間稼働率99.9\%以上を維持する。 \\
拡張性 & 利用者数が10倍になっても構成変更で対応できる。 \\
セキュリティ & 権限のない利用者は個人情報を閲覧できない。 \\
\bottomrule
\end{tabularx}
